\begin{proposition}[Prime ideals of the pointed multiplicative monoid]\label{prop:monoid-spectrum}
Let $M_0$ be the pointed multiplicative monoid underlying $\Z_0$, with new absorbing element $\tau$. Its prime monoid ideals are $\{\tau\}$ and
\[
 P_A=\{\tau,0_\Z\}\cup
 \{n\in\Z\setminus\{0\}:\text{some prime in }A\text{ divides }n\},
\]
one for each subset $A$ of the rational primes, including the empty subset. In particular its monoid spectrum has infinite Krull dimension, and $M_0$ is not Noetherian.
\end{proposition}
\begin{proof}
A monoid ideal contains $\tau$ and is closed under multiplication by arbitrary elements. Its complement must be multiplicatively closed for it to be prime. The set $\{\tau\}$ is prime because products of old ring elements remain old ring elements.

Any larger ideal contains an old integer and hence contains $0_\Z$. A proper ideal contains neither $1$ nor $-1$. By unique factorization of a nonzero integer, primeness then says that it belongs to the ideal exactly when one of its prime divisors does. Thus the ideal is $P_A$, where $A$ is its set of positive prime elements. Conversely the displayed $P_A$ is an ideal and its complement is multiplicatively closed, so it is prime.

Enumerating the rational primes produces arbitrarily long strict chains of the $P_A$, proving infinite dimension. The monoid ideal generated by the first $n$ primes consists of $\tau$, the old zero, and their integer multiples. The next prime belongs to the next such ideal and not the preceding one. These ideals form a strictly ascending infinite chain, proving the failure of the ascending-chain condition.
\end{proof}

The additive semiring spectrum of $\Z_0$ and this multiplicative monoid spectrum are different constructions. Forgetting addition creates all the extra prime ideals $P_A$; it does not preserve the finite-dimensional prime-chain structure of the additive semiring spectrum.
